% Keep styling.sty beside this file. Compile with:
% latexmk -xelatex -synctex=1 template.tex (or, 
% for my purposes, compile to preview with VimtexCompile in neovim)
% if youre using VSC with some built in plugin you'll probably need to add 
% {
%   "latex-workshop.latex.recipe.default": "latexmk (xelatex)",
%   "latex-workshop.latex.autoBuild.run": "onSave",
%   "latex-workshop.view.pdf.viewer": "tab",
%   "latex-workshop.view.pdf.invertMode.enabled": "never"
% }
% also it is necessary to download the right fonts. Just look 
% at error messages to find which ones or something

\documentclass[11pt]{article}
\usepackage{styling} % Dark by default, with the clickable THEME switch.
% \usepackage{graphicx} % Uncomment if adding external images.

\title{Title of the Notes}
\author{Thobias K. H\o ivik}
\date{\today} % Use \date{} to hide the date.

\begin{document}
\maketitle

Just practice problems for Topology, Commutative Algebra, 
and perhaps Machine Learning and Advanced Algorithms.

\section{Exam style stuff}
I have to practice for MAT224 Commutative Algebra 
oral examiniation so these problems are more 
about intuition, allowing less rigor. 

\begin{prob}[Large algebra identification problem]
    Let 
    \[
        R = k[x,y] , \ I = (x^2, xy, y^2), \ A = R/I. 
    \]
    What can you say about $A$? 
    In particular, is it Noetherian or Artinian? 
    Can you give a $k$-basis or normal form for its elements?

    Compute $\sqrt I$. Then describe $V(I) \subseteq k^2$. 
    What information about $I$ is visible from $V(I)$? 

    Is $I$ prime? Is $I$ primary? Give the 
    quickest convincing argument you can. 

    Determine the maximal ideals of $A$. 
    Is $A$ local? 
    If so, identify the maximal ideal 
    $\mathfrak m$, and compute enough powers of 
    $\mathfrak m$ to see what happens. 

    Compute $\Ann_A(\bar x)$ and 
    $\Ann_A(\bar y)$. 
    Are either of these prime ideals? 

    Consider mutliplication by $\bar x$, 
    \[
        \mu_{\bar x} : A \to A, \ a \mapsto \bar x a. 
    \]
    Without being given a sequence, identify 
    $\ker \mu_{\bar x},  \ \im \mu_{\bar x}, \ \coker 
    \mu_{\bar x}$. 

\end{prob}

\begin{proof}[Solution]
    We have $I = (x,y)^2$. We can then say that 
    the ring is Noetherian by using Hilbert's 
    Basis Theorem twice $k \rightarrow k[x] 
    \rightarrow k[x][y] = k[x,y]$ with Noetherian 
    preserved. Then $I \subseteq R$ is 
    Noetherian so $R/I = A$ is Noetherian. 

    Artinian does not hold for $R$ since 
    $(x) \supsetneq (x^2) \supsetneq (x^4) \supsetneq \cdots$ never 
    stabilizes. We have an exact sequence 
    \[
        0 \to (x,y)(x,y) \to R = k[x,y] \to A \to 0.  
    \]
    We have that $R$ is Artinian if 
    and only if $(x,y)^2$ and $A$ are Artinian. 
    We know that $R$ is not Artinian, and 
    $(x,y)^2$ is not Artinian either so we 
    can't conclude anything about $A$ using this. 

    What we can notice is that 
    any polynomials with $x$ degree 
    $\geq 2$ vanish, same for $y$, while 
    terms in both $x$ and $y$ vanish also. Thus 
    all polynomials have only constant terms, terms in $x$ and 
    terms in $y$. 
    Thus $A$ is an $R$-algebra, more specifically a 
    $k$-vector space with basis $\{1,\bar x,\bar y\}$.

    The root $\sqrt I$ is $(x,y)$. 
    We obviously have $(x,y) \subseteq \sqrt I$. 
    For $ax + by \in (x,y)$ we have 
    \[
        (ax + by)^2 = a^2x^2 + (2ab)xy + b^2y^2  
    \] 
    which is in $I$. 
    Now, all generators of $I$ are divisible by by either 
    $x$ or $y$ so $I \subseteq (x,y)$. We have $(x,y)$ prime 
    (and even maximal) so $\sqrt I \subseteq (x,y)$ (since 
    $f^n \in I \subseteq J$ implies $f \in J$ due to 
    primality).

    We do not have that $I$ is prime, since 
    $x^2$ is in $I$, but $x$ is not. It is, however, 
    maximal. We have already established that 
    $(x,y)$ is maximal and $I = (x,y)^2$ so 
    \[
        (x,y)^2 \subseteq I \subseteq (x,y). 
    \]
    Thus $I$ is $(x,y)$-primal. 

    The maximal ideals of $A$ map one to one with the 
    maximal ideals of $R$ containing $I$. 

    As mentioned we have $(x,y) \supseteq (x^2,xy,y^2)$ and 
    $(x,y)$ is maximal so 
    \[
        (\bar x, \bar y) = \frac{(x,y)}{(x^2,xy,y^2)}. 
    \]
    We can use the third isomorphism theorem 
    to prove that $A/(\bar x, \bar y) \cong k$. 

    For uniqueness, 
    suppose $f \in A$. 
    Then $f = a + bx + cy$. 
    The ideal $\mathfrak m = (\bar x,\bar y)$ 
    consists of all elements 
    where the constant term is zero, $a = 0$, 
    so an element $f$ not in 
    $(\bar x,\bar y)$ is just $f = a \in k \setminus \{0\}$, 
    which is 
    invertible since $k$ is a field. 

    As we glimpsed earlier, 
    $\mathfrak m^2 = (\bar x, \bar y)(\bar x,\bar y) 
    = (\bar x^2, \bar x \bar y, \bar y^2) = 0$. 

    We have 
    \[
        \Ann_A(\bar x) = \left\{
            a \in A \bigm| a\bar x = 0. 
        \right\}
    \]
    Equivalently, in $A = R/I$ we have that 
    the annihilator of $\bar x$ are 
    all $a = \alpha + \beta \bar x + \delta \bar y$ such that 
    \[
        a \bar x 
        = \alpha \bar x + \underbrace{\beta \bar x^2}_{0} 
        + \underbrace{\delta \bar y \bar x}_0 = 0. 
    \]
    We have that $\bar x$ is a basiselement of 
    the space, requiring $\alpha = 0$. 
    $\beta, \ \delta$ can be any elements at 
    all so 
    \[
        a = 0 + \beta \bar x + \delta \bar y 
        \in (\bar x, \bar y) = \mathfrak m. 
    \]
    It is not terribly surprising that $\Ann_A(\bar x)$ 
    is contained in $\mathfrak m$, since 
    it is maximal, but actually $\Ann_A(\bar x) = \mathfrak m$. 

    We can see this since, for 
    $a\bar x + b \bar y \in (\bar x, \bar y)$ we have 
    \[
        (a \bar x + b \bar y) \bar x 
        = a \bar x^2 + b \bar y^2 = 0. 
    \]
    Thus $\Ann_A(\bar x) = \mathfrak m$, the maximal ideal. 

    Now, the exact same holds for $\bar y$, for 
    the same reasons, and it is not hard to see. 

    Now for multiplication at $\bar x$ we 
    get a natural short exact sequence 
    \[
        0 \to A
        \overset{\mu_{\bar x}}\to A \overset\pi\to 
        \Ann_A(\bar x) \to 0 
    \]
    
\end{proof}

\end{document}